\section{Congruencias de Hecke, retículos y clases de Selmer}
\label{sec:congruencias-selmer}

La Nota~11 distinguió una representación de Galois semisimple de los
retículos y clases de extensión que puede contener una realización
integral. Ahora hacemos explícito ese paso. El problema de los
denominadores de un proyector de Hecke se reducirá primero a álgebra
lineal sobre un anillo de valoración discreta; después veremos cómo
una representación integral da un cociclo de Galois y qué hipótesis
faltan para convertirlo en una estimación cohomológica útil.

\fuente{A (ejemplar de 55 páginas), §8, pp.~35--42; A93
(ejemplar de 93 páginas), §9, pp.~61--77; B, §5.1--5.2.
Para el desglose de la fuente B se ha consultado además su versión
de 76 páginas del 22 de noviembre de 2026, pp.~54--56,
\url{https://www.ihes.fr/~rzhang/files/GRH\%20note.pdf}.
Esta versión se distingue del ejemplar B de 73 páginas citado en
las notas anteriores. Las proposiciones elementales siguientes
se demuestran aquí; las comparaciones aritméticas globales
se identifican como entradas pendientes. No se ha realizado
un nuevo cotejo de Lean.}

\subsection{Tres problemas que no deben confundirse}

Sean $K$ una extensión finita de $\Q_\ell$, $\mathcal O$ su anillo
de enteros, $\varpi$ un uniformizante, $k=\mathcal O/\varpi\mathcal O$
y $v_\varpi(\varpi)=1$. No confundiremos esta valoración con
$v_\ell(\ell)=1$: si $e=v_\varpi(\ell)$, entonces
$v_\ell(x)=v_\varpi(x)/e$ para $x\in K^\times$.
La longitud $\length_{\mathcal O}$ de un módulo finito de torsión
cuenta sus factores simples $k$.

En el problema aparecen tres objetos distintos:
\begin{enumerate}
\item Un espacio vectorial $V$ sobre $K$, descompuesto en autoespacios
      de operadores de Hecke, y un retículo integral $M\subset V$.
\item Una representación $\rho:G\to\operatorname{GL}(V)$ y un
      retículo $G$-estable, cuya reducción puede contener extensiones
      invisibles en la semisimplificación de $V$.
\item Un grupo cohomológico con condiciones locales, cuya longitud
      puede acotar determinadas clases de extensión \emph{si se
      construye una comparación adecuada}.
\end{enumerate}
La existencia de cualquiera de los tres no proporciona
automáticamente los otros dos. En particular, una congruencia
de autovalores no constituye por sí sola una inyección del
cociente $\mathcal O/\mathfrak c$ asociado a un ideal de
congruencias en un grupo de Selmer.

\subsection{El denominador de una proyección en dimensión dos}

Escribamos $V=Ke_1\oplus Ke_2$ y fijemos $n\ge1$. El retículo
\begin{equation}
 M_n=\mathcal Oe_1+\mathcal O u_n,
 \qquad u_n=\varpi^{-n}e_1+e_2
 \label{eq:cs-reticulo}
\end{equation}
es libre de rango dos; no es la suma de sus intersecciones
con los dos sumandos de $V$. Sea $p:V\to Ke_1$ la proyección
paralela a $Ke_2$.

\begin{proposicion}[Longitud exacta del denominador]
\label{prop:cs-modelo-proyector}
Para el retículo \eqref{eq:cs-reticulo} se cumple
\begin{equation}
 pM_n=\varpi^{-n}\mathcal Oe_1,
 \qquad M_n\cap Ke_1=\mathcal Oe_1,
 \qquad
 \length_{\mathcal O}\frac{pM_n}{M_n\cap Ke_1}=n.
 \label{eq:cs-longitud-exacta}
\end{equation}
\end{proposicion}

\begin{proof}
Como $p(e_1)=e_1$ y $p(u_n)=\varpi^{-n}e_1$,
la primera igualdad es inmediata. Todo elemento de $M_n$ tiene
forma $ae_1+bu_n$ con $a,b\in\mathcal O$; su coordenada $e_2$
es $b$. Estar en $Ke_1$ obliga a $b=0$, de donde resulta la
segunda igualdad. La multiplicación por $\varpi^n$ identifica
el cociente con $\mathcal O/\varpi^n\mathcal O$, que tiene
una filtración de longitud $n$ por sus potencias sucesivas.
\end{proof}

El espacio $V$ está perfectamente descompuesto sobre $K$,
pero $p$ no preserva $M_n$. El defecto no es una cuestión de
normas ni de aproximación: es un cociente de longitud exactamente $n$.

\subsection{La cota general de un operador separador}

\begin{lema}[Proyector y operador integral]
\label{lem:cs-operador}
Sean $V=V_1\oplus V_2$ un espacio vectorial de dimensión finita
sobre $K$, $M$ un retículo completo y $p$ el proyector sobre
$V_1$. Supongamos que $T\in\operatorname{End}_{\mathcal O}(M)$
actúa como el escalar $a\in\mathcal O\setminus\{0\}$ en $V_1$
y como cero en $V_2$. Entonces
\begin{equation}
 apM\subseteq M\cap V_1,
 \qquad
 \length_{\mathcal O}\frac{pM}{M\cap V_1}
 \le (\dim_K V_1)v_\varpi(a).
 \label{eq:cs-operator-bound}
\end{equation}
\end{lema}

\begin{proof}
Para $m\in M$ se tiene $Tm=apm$, y $Tm\in M\cap V_1$.
Por tanto $apM\subseteq M\cap V_1\subseteq pM$.
El módulo $pM$ es un retículo libre de rango
$d=\dim_K V_1$. El cociente del enunciado es cociente de
$pM/apM\simeq(\mathcal O/a\mathcal O)^d$;
su longitud no supera $d\,v_\varpi(a)$.
\end{proof}

En \eqref{eq:cs-reticulo}, el operador $T=\varpi^n p$ es
integral: envía $e_1$ a $\varpi^n e_1$ y $u_n$ a $e_1$.
El lema proporciona $n\le n$, de modo que su cota es
\emph{óptima} en este ejemplo. Este es el mecanismo elemental
de B, Lema 5.1, en §5.1 (B76, p.~54; B73, p.~53);
aquí se ha demostrado de nuevo.

\subsection{Un álgebra de congruencias visible}

Identifiquemos los operadores diagonales sobre $V$ con $K^2$.
Definimos el subanillo
\begin{equation}
 \mathbb T_n=\{(a,b)\in\mathcal O^2:
                a\equiv b\pmod{\varpi^n}\}.
 \label{eq:cs-Tn}
\end{equation}
Los dos homomorfismos $\lambda_1,\lambda_2:\mathbb T_n\to
\mathcal O$ que seleccionan las coordenadas son sistemas de
valores propios congruentes módulo $\varpi^n$.

\begin{proposicion}[El ideal de congruencias del modelo]
\label{prop:cs-ideal}
El anillo $\mathbb T_n$ preserva $M_n$. El ideal
\begin{equation}
 \mathfrak c_1
 =\{\lambda_1(t):t\in\mathbb T_n,
                      \lambda_2(t)=0\}
 \label{eq:cs-ideal}
\end{equation}
es $\varpi^n\mathcal O$. Además,
$(1,0)\notin\mathbb T_n$ mientras que
$(\varpi^n,0)\in\mathbb T_n$. La longitud de
$\mathcal O/\mathfrak c_1$ coincide con el denominador
\eqref{eq:cs-longitud-exacta}.
\end{proposicion}

\begin{proof}
El operador $(a,b)$ envía $e_1$ a $ae_1$ y
\[
 u_n\longmapsto bu_n+\frac{a-b}{\varpi^n}e_1.
\]
Preserva $M_n$ si y solo si $(a-b)/\varpi^n\in\mathcal O$,
que es exactamente la congruencia definitoria de $\mathbb T_n$.
Si $b=0$, esa congruencia equivale a $a\in\varpi^n\mathcal O$.
Las dos afirmaciones de pertenencia a $\mathbb T_n$ se
comprueban sustituyendo $a,b$; por último,
$\length_{\mathcal O}(\mathcal O/\varpi^n\mathcal O)=n$.
\end{proof}

El idempotente racional $(1,0)$ existe en $K^2$, pero
no en el álgebra integral que preserva el retículo.
Multiplicarlo por $\varpi^n$ repara la integralidad al
coste exacto de $n$ unidades de longitud.

\subsection{Una misma representación, dos reducciones}

Supongamos ahora que $G$ es un grupo profinito y que
$\chi_1,\chi_2:G\to\mathcal O^\times$ son caracteres
continuos tales que
\begin{equation}
 \chi_1(g)\equiv\chi_2(g)\pmod{\varpi^n}
 \quad\text{para todo }g\in G.
 \label{eq:cs-character-congruence}
\end{equation}
En $V=Ke_1\oplus Ke_2$ consideremos
$\rho(g)=\operatorname{diag}(\chi_1(g),\chi_2(g))$.
El retículo ordinario $M_0=\mathcal Oe_1\oplus\mathcal Oe_2$
es estable y tiene reducción diagonal. El retículo $M_n$
de \eqref{eq:cs-reticulo} también es estable, pero su reducción
puede no ser diagonalizable.

\begin{proposicion}[Matriz integral y reducción residual]
\label{prop:cs-matriz}
En la base $(e_1,u_n)$ de $M_n$, la representación es
\begin{equation}
 \rho_{M_n}(g)=
 \begin{pmatrix}
 \chi_1(g)&\displaystyle\frac{\chi_1(g)-\chi_2(g)}{\varpi^n}\\
 0&\chi_2(g)
 \end{pmatrix}.
 \label{eq:cs-matriz}
\end{equation}
Su reducción define una extensión
\begin{equation}
 0\longrightarrow k(\bar\chi_1)
 \longrightarrow\bar M_n
 \longrightarrow k(\bar\chi_2)\longrightarrow0.
 \label{eq:cs-extension}
\end{equation}
Los caracteres residuales coinciden, pero la extensión
no tiene por qué escindirse.
\end{proposicion}

\begin{proof}
La igualdad
\[
 \rho(g)u_n
 =\chi_2(g)u_n+
   \frac{\chi_1(g)-\chi_2(g)}{\varpi^n}e_1
\]
da la matriz. La hipótesis \eqref{eq:cs-character-congruence}
hace integral la entrada superior derecha. La primera recta
es estable y el cociente lleva el carácter $\bar\chi_2$.
La igualdad de las diagonales residuales se deduce de la
congruencia, pero no elimina la entrada superior derecha.
\end{proof}

En particular, el hecho de que los polinomios característicos
sean $(X-\chi_1(g))(X-\chi_2(g))$ sobre $K$
no dice si la extensión residual de \eqref{eq:cs-extension}
se escinde. La traza y el determinante no distinguen los
retículos $M_0$ y $M_n$.

\subsection{Del término no diagonal a \texorpdfstring{$H^1$}{H1}}

Dotemos al cuerpo $k$ de la topología discreta y sean
$\alpha,\beta:G\to k^\times$ caracteres continuos y
\begin{equation}
 A(g)=\begin{pmatrix}\alpha(g)&b(g)\\0&\beta(g)\end{pmatrix}
 \label{eq:cs-upper-triangular}
\end{equation}
una representación continua. La identidad $A(gh)=A(g)A(h)$
impone
\begin{equation}
 b(gh)=\alpha(g)b(h)+b(g)\beta(h).
 \label{eq:cs-b-identity}
\end{equation}
El módulo unidimensional
$k(\alpha\beta^{-1})$ es el cuerpo $k$ dotado de la acción
$g\cdot x=\alpha(g)\beta(g)^{-1}x$.

\begin{lema}[Clasificación elemental de la extensión]
\label{lem:cs-cociclo}
Poniendo $c(g)=b(g)/\beta(g)$ se obtiene un cociclo
continuo de grado uno:
\begin{equation}
 c(gh)=c(g)+(\alpha\beta^{-1})(g)c(h).
 \label{eq:cs-cociclo}
\end{equation}
Cambiar un levantamiento de la recta cociente modifica
$c$ en un coborde. La extensión definida por $A$ se
escinde si y solo si su clase
\begin{equation}
 [c]\in H^1\bigl(G,k(\alpha\beta^{-1})\bigr)
 \label{eq:cs-H1}
\end{equation}
es cero.
\end{lema}

\begin{proof}
Dividimos \eqref{eq:cs-b-identity} entre
$\beta(gh)=\beta(g)\beta(h)$, y resulta
\eqref{eq:cs-cociclo}. Cambiar el segundo vector de la base
por $e_2+t e_1$ reemplaza $b(g)$ por
$b(g)+t(\alpha(g)-\beta(g))$ y, por tanto, $c(g)$ por
\[
 c(g)+t\bigl((\alpha\beta^{-1})(g)-1\bigr).
\]
El segundo sumando es un coborde. Una recta complementaria
$G$-estable equivale a elegir $t$ que anule la entrada
superior derecha para todos los $g$; esa condición equivale
a que el cociclo sea un coborde.
\end{proof}

Es importante el orden de los caracteres:
$\operatorname{Hom}_k(k(\beta),k(\alpha))$ tiene acción
$\alpha\beta^{-1}$, \emph{no} $\beta\alpha^{-1}$.
Cuando $\alpha=\beta$, la acción es trivial y los
cobordes son cero. En ese caso, cualquier cociclo no nulo
produce una extensión no escindida.

\subsection{Ejemplo completo sobre \texorpdfstring{$\Z_\ell$}{Z-ell}}

Tomemos $K=\Q_\ell$, $\mathcal O=\Z_\ell$,
$\varpi=\ell$ y $G=(\Z_\ell,+)$.
Sea $\chi_1(z)=1$ y definamos
$\chi_2(z)=(1+\ell)^z$ por el desarrollo binomial
$\ell$-ádico. Es un carácter continuo con valores en
$\Z_\ell^\times$, y
\begin{equation}
 (1+\ell)^z\equiv1+\ell z\pmod{\ell^2}.
 \label{eq:cs-binomial}
\end{equation}
La congruencia se comprueba primero para enteros
no negativos y se extiende a $\Z_\ell$ por continuidad.
El desarrollo binomial converge uniformemente en $\Z_\ell$,
pues $\binom{z}{m}\in\Z_\ell$ y $\ell^m$ tiende a cero.
La identidad multiplicativa también se extiende por continuidad.
El argumento incluye el primo $\ell=2$.

\begin{proposicion}[Extensión no trivial con semisimplificación fija]
\label{prop:cs-Zell}
Para $n=1$, la reducción de $M_1$ es
\begin{equation}
 \bar\rho_{M_1}(z)
 =\begin{pmatrix}1&-\bar z\\0&1\end{pmatrix}.
 \label{eq:cs-unipotente}
\end{equation}
La sucesión \eqref{eq:cs-extension} no se escinde,
mientras que la reducción de $M_0$ sí se escinde.
Ambas provienen de la misma representación semisimple
$1\oplus\chi_2$ sobre $\Q_\ell$.
\end{proposicion}

\begin{proof}
La entrada superior derecha de \eqref{eq:cs-matriz}
es $(1-(1+\ell)^z)/\ell$, cuya reducción es $-\bar z$
por \eqref{eq:cs-binomial}. Como la acción residual
sobre los coeficientes es trivial, el cociclo $z\mapsto-\bar z$
representa un elemento de
$H^1_{\mathrm{cont}}(\Z_\ell,\mathbb F_\ell)$.
No es cero porque vale $-1$ en $z=1$.
Con acción trivial todos los cobordes son nulos; por el
Lema~\ref{lem:cs-cociclo} la extensión no se escinde.
La reducción del retículo $M_0$ es la suma directa
de las dos rectas triviales. La representación sobre
$\Q_\ell$ no ha cambiado.
\end{proof}

Esta diferencia es el motivo para no inferir una clase
integral de la mera existencia de la representación
semisimple de la Nota~11. La congruencia de trazas no garantiza
que la reducción de un retículo estable elegido arbitrariamente
sea no escindida: $M_0$ tiene reducción escindida, mientras que
$M_1$ conserva la clase residual no nula.

\paragraph{La clase residual y el cambio de coeficientes.}
La no escisión residual tampoco garantiza que la clase sobreviva
en el módulo divisible que se obtiene de los mismos caracteres.
Volvamos al modelo general y pongamos
\[
 \eta=\chi_1\chi_2^{-1},\qquad
 c_n(g)=\frac{\eta(g)-1}{\varpi^n}\in\mathcal O.
\]
Este es el cociclo normalizado de \eqref{eq:cs-matriz}:
$c_n=b/\chi_2$. Su reducción $\bar c_n$ clasifica la extensión
residual. Sea $A=(K/\mathcal O)(\eta)$, con la acción
$g\cdot x=\eta(g)x$. La aplicación
\[
 j:k(\bar\eta)\hookrightarrow A,\qquad
 \bar x\longmapsto x/\varpi+\mathcal O
\]
es una inyección bien definida y $G$-equivariante: cambiar
el representante $x$ por $x+\varpi y$ no altera su imagen,
y la reducción de la acción en $A[\varpi]$ es $\bar\eta$.
Sin embargo,
\[
 \begin{aligned}
 j\bigl(\bar c_n(g)\bigr)
 &=\frac{\eta(g)-1}{\varpi^{n+1}}+\mathcal O\\
 &=g\cdot\bigl(\varpi^{-(n+1)}+\mathcal O\bigr)
   -\bigl(\varpi^{-(n+1)}+\mathcal O\bigr).
 \end{aligned}
\]
Por tanto $j_*[\bar c_n]=0$ en $H^1(G,A)$: la imagen
es un coborde. En el ejemplo anterior,
$\eta(z)=(1+\ell)^{-z}$ y $n=1$, de modo que la clase
residual $z\mapsto-\bar z$, que es no nula, se convierte
en el coborde de $\ell^{-2}+\Z_\ell$.

Una inyección de coeficientes no tiene por qué inducir una
inyección en $H^1$. Aquí la sucesión
$0\to A[\varpi]\to A\xrightarrow{\varpi}A\to0$
lo explica: $\varpi^{-n}+\mathcal O$ es invariante y su
imagen por la aplicación de conexión es $[\bar c_n]$,
tras identificar $A[\varpi]$ con $k(\bar\eta)$ mediante $j$.
El cálculo se refiere a la acción levantada $\eta$;
no autoriza a cambiarla por otra acción en $K/\mathcal O$.
Para usar coeficientes divisibles en Selmer habrá que comprobar,
además de las condiciones locales, que las clases pertinentes
sobreviven a la comparación de coeficientes.

\subsection{Restricciones locales y condición de Selmer}

Sea ahora $E$ un cuerpo de números, $S$ un conjunto finito
de lugares finitos y $G_{E,S}$ el grupo de Galois de la
máxima extensión de $E$ no ramificada fuera de $S$ y de
los lugares infinitos. En las aplicaciones $\ell$-ádicas
incluimos en $S$ los lugares sobre $\ell$ y todos los de
ramificación de los coeficientes. Fijemos un $\mathcal O$-módulo
de torsión discreto $A$ con acción continua y $\mathcal O$-lineal
de $G_{E,S}$ y, para cada $w\in S$, un $\mathcal O$-submódulo
prescrito $L_w\subseteq H^1(E_w,A)$.

\begin{definicion}[Grupo de Selmer con condiciones prescritas]
\label{def:cs-selmer}
El grupo correspondiente es
\begin{equation}
 H^1_{\mathcal L}(E,A)
 =\ker\!\left(
 H^1(G_{E,S},A)\longrightarrow
 \prod_{w\in S}\frac{H^1(E_w,A)}{L_w}
 \right).
 \label{eq:cs-selmer}
\end{equation}
En los lugares fuera de $S$ la no ramificación ya está
incorporada en el grupo global. La notación no supone que
las condiciones $L_w$ sean iguales para todos los lugares.
Esta definición no impone condiciones en los lugares reales;
si se desean, deben añadirse. En el caso CM posterior no hay
lugares reales.
\end{definicion}

Si $T\subset V$ es un retículo estable en una representación
$\ell$-ádica, se utiliza con frecuencia
$A=V/T$, módulo de torsión que no es en general finito.
En lugares sobre $\ell$, una elección aritmética importante
es la condición finita de Bloch--Kato: la imagen de
$H^1_f(E_w,V)$ en $H^1(E_w,A)$.
Fuera de $\ell$ y de los lugares permitidos se exige
no ramificación; en ciertos lugares permitidos puede
admitirse todo $H^1(E_w,A)$. Estas son condiciones distintas.

\begin{lema}[Escisión tras restringir a un grupo local]
\label{lem:cs-local-split}
Una extensión de caracteres sobre $G$, clasificada por
$[c]\in H^1(G,k(\alpha\beta^{-1}))$, se escinde al
restringirse a un subgrupo cerrado $D\subseteq G$ si
y solo si
$\operatorname{res}_D[c]=0$ en
$H^1(D,k(\alpha\beta^{-1}))$.
\end{lema}

\begin{proof}
La matriz de la representación restringida tiene la
misma forma \eqref{eq:cs-upper-triangular}, ahora solo
para $g\in D$. El argumento de cambio de levantamiento
del Lema~\ref{lem:cs-cociclo} se aplica a $D$, de modo
que existe un complemento $D$-estable exactamente cuando
la clase restringida es un coborde.
\end{proof}

La condición $\operatorname{res}_D[c]=0$ es más fuerte
que pedir que la clase pertenezca a un subgrupo local
admisible $L_w$ que no sea cero. No reemplazaremos la
condición finita de Bloch--Kato por una afirmación de escisión,
ni la no ramificación por trivialidad local.

\subsection{Una clase puede ser globalmente no trivial y localmente nula}

En el ejemplo anterior, fijemos el subgrupo cerrado
$D=\ell\Z_\ell\subseteq G=\Z_\ell$. El cociclo
$c(z)=-\bar z$ no es cero en $H^1(G,\mathbb F_\ell)$, pero
se anula sobre $D$. Por el Lema~\ref{lem:cs-local-split},
la extensión no se escinde sobre $G$ y sí se escinde sobre
$D$. Esto modela por qué una condición local de anulación
no implica por sí sola la anulación global.

En un cuerpo de números no se puede elegir arbitrariamente
los grupos $D_w$: provienen de grupos de descomposición,
y se debe comprobar que la clase cumple \emph{todas} las
condiciones locales. El ejemplo es un modelo cohomológico,
no una construcción de una clase global de Selmer para un
cuerpo de números concreto.

\subsection{Del ideal de congruencias a una cota de Selmer}

La longitud del denominador de proyección está controlada
por un ideal de congruencias, como en
\eqref{eq:cs-operator-bound} y \eqref{eq:cs-ideal}.
El paso verdaderamente aritmético consiste en relacionar
ese ideal con clases de extensión sujetas a condiciones
locales. La siguiente proposición enuncia una hipótesis
suficiente para obtener esa comparación de longitudes.

\begin{proposicion}[Puente condicional de longitudes]
\label{prop:cs-puente}
Supongamos que $\mathfrak c=(\varpi^n)\ne0$ es un
ideal de congruencias, $S_\mathcal L$ un $\mathcal O$-módulo
de longitud finita y existe una inyección $\mathcal O$-lineal
\begin{equation}
 \iota:\mathcal O/\mathfrak c\hookrightarrow S_\mathcal L.
 \label{eq:cs-inyeccion}
\end{equation}
Entonces $n\le\length_{\mathcal O}S_\mathcal L$.
Si, además, el operador integral del
Lema~\ref{lem:cs-operador} tiene escalar $a$ con
$v_\varpi(a)=n$ y $d=\dim_KV_1$, entonces
\begin{equation}
 \length_{\mathcal O}\frac{pM}{M\cap V_1}
 \le d\,\length_{\mathcal O} S_\mathcal L.
 \label{eq:cs-selmer-bound}
\end{equation}
\end{proposicion}

\begin{proof}
Una inyección de módulos de longitud finita no
puede aumentar la longitud:
$\length_{\mathcal O}(\mathcal O/\mathfrak c)
=n\le\length_{\mathcal O}S_\mathcal L$.
La segunda afirmación resulta al concatenar esta
desigualdad con \eqref{eq:cs-operator-bound}.
\end{proof}

Más generalmente, sea $f:\mathcal O/\mathfrak c\to S_\mathcal L$
un mapa $\mathcal O$-lineal con
$\length_{\mathcal O}\ker f\le\kappa$. La aditividad de la
longitud en la sucesión exacta de núcleo e imagen da
\[
 n=\length_{\mathcal O}\ker f+
   \length_{\mathcal O}\operatorname{im}f
 \le\kappa+\length_{\mathcal O}S_\mathcal L.
\]
Bajo la misma hipótesis sobre el escalar $a$, se obtiene
\[
 \length_{\mathcal O}\frac{pM}{M\cap V_1}
 \le d\bigl(\kappa+\length_{\mathcal O}S_\mathcal L\bigr).
\]
La inyección es el caso $\kappa=0$. La desigualdad de
longitudes por sí sola no implica que exista tal inyección,
como muestra el Ejercicio~\ref{ej:cs-longitudes}.

La proposición \emph{no afirma} que la inyección
\eqref{eq:cs-inyeccion} exista para las álgebras de Hecke
del argumento de GRH. Construir un mapa con suficiente
control de su núcleo, del retículo, de las condiciones locales
y de las normalizaciones es parte de la entrada profunda.
Para varios sistemas de autovalores, el mismo razonamiento
se aplica a cada cociente de congruencias si se dispone
de los mapas y cotas correspondientes; las dimensiones
de los autoespacios multiplican las pérdidas.

\subsection{Obstrucciones al levantamiento: el lugar de \texorpdfstring{$H^2$}{H2}}

El primer grupo cohomológico describe extensiones;
el segundo aparece al intentar elevar cociclos a un nivel
integral superior. Este fenómeno se observa sin introducir
deformaciones automorfas.

\begin{proposicion}[Bockstein para módulos finitos]
\label{prop:cs-bockstein}
Sea $T$ un $\mathcal O$-módulo libre de rango finito
con acción continua y $\mathcal O$-lineal de $G$.
Para $r\ge1$ existe una sucesión exacta de módulos
finitos discretos de $G$:
\begin{equation}
 0\longrightarrow T/\varpi^r T
 \xrightarrow{\ \times\varpi\ }
 T/\varpi^{r+1}T
 \longrightarrow T/\varpi T\longrightarrow0.
 \label{eq:cs-exacta-bockstein}
\end{equation}
Una clase $x\in H^1(G,T/\varpi T)$ se eleva a
$H^1(G,T/\varpi^{r+1}T)$ si y solo si su obstrucción
\begin{equation}
 \delta_r(x)\in H^2(G,T/\varpi^rT)
 \label{eq:cs-bockstein}
\end{equation}
es cero.
\end{proposicion}

\begin{proof}
La multiplicación por $\varpi$ es inyectiva
desde $T/\varpi^rT$ a $T/\varpi^{r+1}T$ porque
$T$ no tiene torsión; su imagen es el núcleo de
la reducción módulo $\varpi$.
La sucesión larga de cohomología continua contiene
\[
 H^1(G,T/\varpi^{r+1}T)
 \longrightarrow H^1(G,T/\varpi T)
 \xrightarrow{\delta_r} H^2(G,T/\varpi^rT).
\]
La exactitud identifica la imagen de la primera flecha
con el núcleo de $\delta_r$, que demuestra la equivalencia.
\end{proof}

En términos de cociclos, se levanta provisionalmente
una función a nivel $\varpi^{r+1}$. El defecto de satisfacer
la identidad de cociclo toma valores en el núcleo
$T/\varpi^rT$ y produce un $2$-cociclo. Su clase se
anula exactamente cuando puede corregirse el levantamiento.
El grupo de esta obstrucción tiene coeficientes finitos
$T/\varpi^rT$. Su aparición no lo identifica con el grupo
$H^2(G_{E,S},\Z_\ell(j+1))$ que interviene en la comparación
dual de la Nota~8. Esa identificación exigiría fijar
coeficientes, condiciones locales y mapas de comparación;
las cotas aritméticas de la Nota~8 requieren además sus
propias hipótesis. La proposición anterior demuestra el
criterio formal de levantamiento.

\subsection{Normalizaciones y torcidos de Tate}

La descomposición que la fuente B utiliza en §5.2
(B76, p.~55; B73, p.~54) para un
levantamiento theta tiene, después de fijar las entradas,
la forma semisimple
\begin{equation}
 R_\lambda^{\mathrm{ss}}
 \simeq\rho_\lambda\oplus
 \bigoplus_{c=1}^{h}\xi_c,
 \qquad
 \xi_c=\xi_1\chi_\ell^{-(c-1)}.
 \label{eq:cs-twists}
\end{equation}
Tras torcer por $\xi_1^{-1}$, los sumandos de dimensión
uno son
\[
 1,\ \chi_\ell^{-1},\ \ldots,\
 \chi_\ell^{-(h-1)}.
\]
Para una extensión con subobjeto de carácter $\xi_c$ y
cociente de carácter $\xi_d$, el módulo de cociclos es
$\operatorname{Hom}(\xi_d,\xi_c)$. Si $c>d$, su carácter es
\[
 \xi_c\xi_d^{-1}=\chi_\ell^{-(c-d)}.
\]
De esta orientación proceden los torcidos negativos $(-j)$
con $1\le j<h$; la orientación inversa produce $+j$.
La selección de los bloques y de sus condiciones locales
forma parte del enlace aritmético.

La fuente organiza términos con coeficientes
$\Q_\ell/\Z_\ell(-j)$ y otro término con coeficientes
$A_\pi=(T_\pi\otimes_{\mathcal O}K)/T_\pi$, donde
$T_\pi$ es un retículo $G_E$-estable en
$\rho_\lambda\otimes\xi_1^{-1}$.
No son necesariamente grupos de Selmer con las mismas
condiciones locales: el término de dimensión mayor
usa la condición finita en los lugares sobre $\ell$,
mientras que el grupo de torsión unidimensional de la
Nota~8 se organiza como un núcleo de restricción local.

Si $E$ es CM y $S_\ell$ es el conjunto de lugares sobre
$\ell$, el segundo tipo se escribe
\begin{equation}
 S_\ell(E,j)=\ker\!\left(
 H^1\bigl(G_{E,S_\ell},\Q_\ell/\Z_\ell(-j)\bigr)
 \longrightarrow
 \prod_{w\mid\ell}
 H^1\bigl(E_w,\Q_\ell/\Z_\ell(-j)\bigr)
 \right),\quad j\ge1.
 \label{eq:cs-Sell}
\end{equation}
Este grupo es $\operatorname{Sha}_{\ell,j}(E)$ en la
notación de la Nota~8; aquí $S_\ell(E,j)$ conserva la
denominación de B, §5.2 (B76, pp.~55--56; B73, pp.~54--55).
Con Frobenius geométrico, el giro $(-j)$ actúa
fuera de $\ell$ por $q_w^j$, de acuerdo con
la convención de las Notas 8 y 11.
Los cuantificadores sobre $E$, $\ell$, $j$, los retículos
y las extensiones auxiliares son parte de cualquier
estimación \emph{uniforme}: no quedan fijados por
una igualdad formal de grupos.

\subsection{Qué ha quedado demostrado y qué falta}

Los cálculos de esta nota prueban que (i) un proyector
racional puede introducir un denominador de longitud
arbitraria, (ii) el ideal de congruencias mide exactamente
ese defecto en el modelo $\mathbb T_n$, (iii) una elección
de retículo puede retener una clase de $H^1$ no escindida
que otra elección elimina, (iv) esa clase residual puede
anularse al pasar al módulo divisible con la acción levantada,
y (v) el levantamiento a un nivel integral superior puede
presentar una obstrucción en $H^2$.

Para obtener una cota de denominadores en el argumento
aritmético completo todavía es necesario, entre otros pasos,
construir las representaciones y retículos compatibles
con los sistemas de Hecke; producir la comparación entre
congruencias y clases de extensión con el control suficiente
sobre núcleo, cambio de coeficientes y multiplicidades;
demostrar que las clases
satisfacen las condiciones locales requeridas; y obtener
cotas uniformes de los grupos cohomológicos resultantes.
En particular, ni \eqref{eq:cs-selmer-bound} ni la
Proposición~\ref{prop:cs-Zell} sustituyen los teoremas de
deformación, automorfía, dualidad y control de torre
que intervienen en A, A93 y B.

\subsection{Ejercicios resueltos}

\begin{ejercicio}\label{ej:cs-reticulo}
Sea $n=3$, $\mathcal O=\Z_5$ y
$M=\mathcal Oe_1+\mathcal O(5^{-3}e_1+e_2)$.
Calcule el cociente de proyección y el menor exponente
$m\ge0$ para el cual $5^m p$ preserva $M$.
\end{ejercicio}

\paragraph{Solución del Ejercicio~\ref{ej:cs-reticulo}.}
El cociente es $5^{-3}\Z_5/\Z_5\simeq\Z_5/5^3\Z_5$,
por lo que su longitud es tres. La imagen de
$5^mp(5^{-3}e_1+e_2)$ es $5^{m-3}e_1$.
Está en $M\cap Ke_1=\Z_5e_1$ si y solo si $m\ge3$.
El menor exponente es $m=3$.

\begin{ejercicio}\label{ej:cs-ideal}
Para $\mathbb T_2\subset\Z_3^2$, determine
$\mathfrak c_1$ y compruebe directamente
que $(3,0)$ no preserva $M_2$ mientras que
$(9,0)$ sí lo hace.
\end{ejercicio}

\paragraph{Solución del Ejercicio~\ref{ej:cs-ideal}.}
La condición es $a\equiv b\pmod9$; con $b=0$
resulta $\mathfrak c_1=9\Z_3$.
El operador $(3,0)$ envía
$u_2=3^{-2}e_1+e_2$ a $3^{-1}e_1$, que no está
en $M_2\cap Ke_1=\Z_3e_1$.
El operador $(9,0)$ lo envía a $e_1\in M_2$.

\begin{ejercicio}\label{ej:cs-cociclo}
Si $A(g)$ es la matriz \eqref{eq:cs-upper-triangular},
verifique la identidad de cociclo para
$c=b/\beta$ y determine la clase de coeficientes
cuando $\alpha=\beta$.
\end{ejercicio}

\paragraph{Solución del Ejercicio~\ref{ej:cs-cociclo}.}
La multiplicación de matrices da
$b(gh)=\alpha(g)b(h)+b(g)\beta(h)$.
Dividir por $\beta(g)\beta(h)$ produce
$c(gh)=(\alpha\beta^{-1})(g)c(h)+c(g)$.
Si $\alpha=\beta$, el cociclo toma valores en
$k$ con acción trivial; es un homomorfismo
continuo de $G$ a su grupo aditivo. Los cobordes
son cero.

\begin{ejercicio}\label{ej:cs-local}
En $G=\Z_\ell$, con $D=\ell\Z_\ell$ y
$c(z)=-z\bmod\ell$, determine si la clase
es trivial en $G$ y si lo es en $D$.
\end{ejercicio}

\paragraph{Solución del Ejercicio~\ref{ej:cs-local}.}
La acción es trivial, de modo que los cobordes
son cero. Como $c(1)=-1\ne0$, la clase global
no es nula. Para $z\in D$ se tiene $z=\ell t$
y $c(z)=0$, así que la restricción es nula.
La extensión es escindida sobre $D$ pero no
sobre $G$.

\begin{ejercicio}\label{ej:cs-longitudes}
Suponga que existe una inyección
$\mathcal O/\varpi^n\mathcal O\hookrightarrow
S$ con
$S\simeq\mathcal O/\varpi^2\mathcal O\oplus
\mathcal O/\varpi^3\mathcal O$.
Compare la cota de longitud con la cota que se
obtiene del exponente respecto de $\varpi$ de $S$.
\end{ejercicio}

\paragraph{Solución del Ejercicio~\ref{ej:cs-longitudes}.}
La longitud de $S$ es $2+3=5$, y la inyección
implica inicialmente $n\le5$. Pero todo elemento
de $S$ es anulado por $\varpi^3$. Por la inyectividad,
la imagen $x$ de $1$ tiene anulador
$\operatorname{Ann}_{\mathcal O}(x)=(\varpi^n)$.
Como $\varpi^3x=0$, se obtiene
$\varpi^3\in(\varpi^n)$, es decir, $n\le3$.
La cota por el exponente respecto de $\varpi$ es aquí más fuerte.
La existencia de una inyección impone más que
la desigualdad de longitudes aislada.
